\section{Contexto y motivación}

Las personas con ceguera o baja visión enfrentan dificultades significativas al desplazarse por entornos desconocidos o dinámicos.
Las herramientas de asistencia tradicionales, como el bastón blanco o los perros guía, presentan limitaciones en alcance,
resolución espacial y capacidad de anticipación frente a obstáculos no detectables a corta distancia.
En particular, estos métodos no advierten eficazmente objetos a altura de cabeza (ramas, cables, señales colgantes u otros
elementos sobresalientes) cuya detección requiere cobertura vertical que los sensores puntuales no pueden proveer.

En los últimos años han surgido dispositivos electrónicos de asistencia que incorporan sensores activos
\cite{Novich2015, Santos2023, Petsiuk2019UltrasoundNavigation}.
Sin embargo, gran parte de esas soluciones se basa en sensores puntuales de corto alcance, lo que restringe la cantidad
y calidad de información espacial disponible para el usuario.
En contraste, los sensores de profundidad por visión estéreo activa, como la familia Intel RealSense \cite{RS-D400},
permiten obtener un mapa denso del entorno con cobertura bidimensional del campo visual, manteniendo al mismo tiempo portabilidad y baja latencia.

El presente trabajo se enmarca en el principio de neuroplasticidad sensorial, según el cual el cerebro puede reorganizarse
para explotar de manera más eficiente los sentidos restantes cuando la visión está reducida o ausente.
Sobre esa base, el proyecto global del que forma parte esta tesis busca diseñar un sistema portátil capaz de traducir
información espacial visual en estímulos hápticos interpretables para la navegación asistida.

\section{Planteo del problema}

El problema de ingeniería abordado en esta tesis consiste en diseñar un sistema embebido capaz de:
\begin{enumerate}
  \item adquirir información tridimensional del entorno mediante una cámara RGB-D con IMU integrada
  \item procesarla en tiempo real sobre hardware de cómputo limitado
  \item generar una representación compacta, estable y de baja latencia del espacio inmediato frente al usuario, adecuada para ser consumida por una futura interfaz háptica.
\end{enumerate}

Las restricciones determinantes del diseño son: portabilidad (el dispositivo se porta en la cabeza del usuario), cómputo embebido limitado (plataforma del orden de una Raspberry Pi 5 \cite{RaspberryPi}), latencia extremo a extremo reducida y estabilidad temporal de la representación frente a oclusiones cambios de orientación del sensor.

\section{Objetivos}

\subsection{Objetivo del proyecto}

El objetivo final del proyecto es desarrollar un sistema portátil de asistencia sensorial que permita a personas con ceguera
o baja visión percibir su entorno mediante estímulos táctiles, mejorando su seguridad, autonomía
y capacidad de desplazamiento independiente.
El proyecto se estructura en dos etapas complementarias: (i)~captación del entorno y mapeo local (objeto de esta tesis) y
(ii)~traducción de dicha representación a una interfaz háptica, que constituirá un trabajo separado.

\subsection{Objetivo de la tesis}

Esta tesis aborda exclusivamente la primera etapa: la captación del entorno y la generación de un mapeo local compacto
y estable en tiempo real a partir de percepción RGB-D.

\subsection{Objetivo técnico}

El objetivo técnico es desarrollar e implementar en ROS2 \cite{ROS} un módulo de percepción que, utilizando la cámara
Intel RealSense D435i \cite{RS-D400}, genere en tiempo real una representación compacta del entorno a partir de información
de profundidad, optimizada para su uso como entrada de una futura interfaz háptica de asistencia a la movilidad.

Específicamente, el módulo deberá:
\begin{itemize}
  \item procesar imágenes de profundidad \texttt{16UC1} mediante filtrado y reducción de dimensionalidad,
        produciendo como salida una grilla estructurada de baja resolución (por ejemplo, $10\times5$ celdas) con estadísticas
        por celda (media, mínimo, conteo de puntos válidos);
  \item publicar resultados en tópicos ROS2 estandarizados y proveer una interfaz serial para la transmisión eficiente
        de la grilla compacta hacia un microcontrolador externo;
  \item implementar una capa de mapeo local de corto alcance basada en grillas de ocupación log-odds \cite{Elfes1989,Thrun2005},
        con memoria temporal de obstáculos y mecanismo de desplazamiento egocéntrico (\emph{scroll});
  \item incorporar información de la IMU integrada del sensor para estimar la orientación y compensar cambios de pose
        durante la actualización del mapa local.
\end{itemize}

Los parámetros numéricos y umbrales de decisión se refinarán mediante una fase de calibración y pruebas experimentales.

\section{Alcance}

El alcance de esta tesis comprende:
\begin{itemize}
  \item el desarrollo del módulo de percepción y la generación de la representación compacta instantánea de profundidad
  \item la capa de mapeo local de corto alcance con memoria temporal e integración de información inercial
  \item la publicación de resultados en ROS2 y la interfaz serial hacia módulos externos
  \item la evaluación experimental del sistema mediante métricas cuantitativas de precisión, latencia y consumo de recursos
\end{itemize}

Quedan \emph{fuera} del alcance de esta tesis el diseño, la implementación y la integración de la interfaz háptica
y sus actuadores, así como la validación con usuarios finales.
La plataforma embebida objetivo se seleccionará a partir de las mediciones de rendimiento (CPU, memoria y latencia)
obtenidas durante el prototipado; la referencia inicial es una Raspberry Pi 5, con una plataforma alternativa
(NVIDIA Jetson Nano \cite{JetsonNano}) en caso de que los requisitos temporales no sean satisfechos.

\section{Organización del documento}

El presente documento se organiza de la siguiente manera.
El Capítulo~\ref{chap:marco-teorico} desarrolla los fundamentos teóricos del sistema:
percepción de profundidad por estereoscopía activa, modelo de cámara \emph{pinhole}, reducción espacial
mediante partición en regiones de interés, grillas de ocupación log-odds, algoritmo de Bresenham,
estimación de orientación con IMU y protocolo de comunicación serie.
El Capítulo~\ref{chap:estado-arte} revisa el estado del arte en dispositivos de asistencia sensorial con
retroalimentación háptica y en técnicas de mapeo local y evitación de obstáculos en robótica móvil,
posicionando la contribución de esta tesis respecto de las soluciones existentes.
El Capítulo~\ref{chap:arquitectura} describe la arquitectura modular del sistema implementada en ROS2,
con sus dos \emph{pipelines} paralelos —captación compacta directa y mapeo local— y la topología de tópicos.
El Capítulo~\ref{chap:metodo} detalla el diseño del \emph{pipeline} de percepción: preprocesamiento,
agregación y reducción de dimensionalidad, representación instantánea, mapa local con grilla de ocupación
y mecanismo de \emph{scroll} egocéntrico, compensación gravitacional y estimación de altura del suelo.
El Capítulo~\ref{chap:implementacion} expone la implementación concreta en C++ y Python sobre ROS2,
describiendo los nodos principales (\texttt{depth\_to\_matrix\_node}, \texttt{local\_mapper\_node},
\texttt{haptic\_grid\_generator}, \texttt{hardware\_manager}) y las interfaces definidas.
El Capítulo~\ref{chap:experimentos} presenta la metodología experimental, los escenarios de prueba y
las métricas de evaluación utilizadas: error RMS de profundidad, latencia \emph{end-to-end},
precisión y \textit{recall} de detección de obstáculos, estabilidad temporal por celda y consumo de recursos.
El Capítulo~\ref{chap:resultados} expone los resultados empíricos obtenidos y su análisis cuantitativo.
Finalmente, los Capítulos~\ref{chap:discusion}, \ref{chap:conclusiones} y~\ref{chap:futuro} discuten
las limitaciones del enfoque, formulan las conclusiones del trabajo y proponen líneas de trabajo futuro,
respectivamente.
